% ============================================================================
% SECTION 6: ACTIVE-EIGENSPACE CURVATURE
% ============================================================================
\section{Active-Eigenspace Curvature: Computation and Status}\label{sec:curvature}

\subsection{What curvature is being computed}\label{subsec:curvature-status}

The regular Walsh exponential family is dually flat with respect to the e- and
m-connections, but its Levi--Civita curvature is not automatically zero.  At
exact key distributions $P_k$, however, the ordinary Riemannian tensor of the
full-support family is not directly defined because $P_k$ lies on the boundary.
This section therefore uses the following restricted object.

\begin{definition}[Active covariance eigenspace]\label{def:active-eigenspace}
Let $\mathcal I^{\partial}(k)$ be the boundary covariance matrix in
Definition~\ref{def:boundary-fisher}.  Its active eigenspace is the span of the
eigenvectors with strictly positive eigenvalues.  Curvature computations below
are performed after projecting the covariance and third cumulant tensors to this
active eigenspace.
\end{definition}

The resulting values are curvature diagnostics of the limiting regularized
geometry obtained from the active covariance eigenspace.  The finite orbit
$\{P_k\}$ itself is discrete and is used here through its induced boundary
covariance and cumulant tensors.

\subsection{Universal rank and eigenvalues of the boundary covariance}\label{subsec:rank-eigen}

\begin{theorem}[Rank and active eigenvalues]\label{thm:active-rank}
For any bijection $S:\B{s}\to\B{s}$, the boundary covariance matrix
$\mathcal I^{\partial}(k)$ has rank $2^s-1$.  Its nonzero eigenvalues are all
$2^s$ when the feature coordinates are the unnormalised nonconstant Walsh
characters used in Definition~\ref{def:walsh-exp}.
\end{theorem}

\begin{proof}
Let $N=2^s$ and index rows by $a\in\B{s}$.  Let $X$ be the $N\times(2^{2s}-1)$
feature matrix with entries
$X_{a,(\alpha,\beta)}=(-1)^{\alpha\cdot a\oplus\beta\cdot S(a\oplus k)}$ for
$(\alpha,
\beta)\neq(0,0)$.  The boundary covariance is
$N^{-1}X^\top H X$, where $H=I_N-N^{-1}\mathbf 1\mathbf 1^\top$ centers over the
support graph.  The nonzero eigenvalues of $N^{-1}X^\top H X$ equal those of
$N^{-1}HXX^\top H$.  For $a,a'\in\B{s}$,
\[
 (XX^\top)_{a,a'}=
 \sum_{(\alpha,\beta)\neq(0,0)}
 (-1)^{\alpha\cdot(a\oplus a')\oplus
 \beta\cdot(S(a\oplus k)\oplus S(a'\oplus k))}
 =N^2\mathbf 1[a=a']-1,
\]
because $S$ is injective.  Hence
$N^{-1}HXX^\top H=N^{-1}H(N^2I_N-\mathbf 1\mathbf 1^\top)H=NH$.
The matrix $H$ has rank $N-1$ and eigenvalue $1$ on the centered subspace.
Therefore the active eigenvalues are all $N=2^s$.
\end{proof}

\subsection{Third cumulants and curvature protocol}\label{subsec:third-cum}

Let
\begin{equation}\label{eq:third-cumulant}
  C^{\partial}_{ijk}(k)=
  \mathbb E_{P_k}\left[(T_i-\eta_i)(T_j-\eta_j)(T_k-\eta_k)\right]
\end{equation}
be the boundary third cumulant tensor.  Since products of Walsh characters are
again Walsh characters, every entry can be computed exactly from the WHT of
$S$ and the key-translation rule.  If $U$ is an orthonormal basis of the active
eigenspace and $\lambda_a=2^s$ are the active covariance eigenvalues, the
projected tensor is
\begin{equation}\label{eq:projected-cumulant}
  \widetilde C_{abc}=\sum_{i,j,\ell}U_{ia}U_{jb}U_{\ell c}C^{\partial}_{ij\ell}.
\end{equation}
The Levi--Civita curvature diagnostic used in the experiments is then computed
with the standard dually-flat formula
\begin{equation}\label{eq:curvature-diagnostic}
  R_{abab}=\frac14\sum_c
  \frac{\widetilde C_{aac}\widetilde C_{bbc}-\widetilde C_{abc}^2}{\lambda_c},
  \qquad
  K_{ab}=\frac{R_{abab}}{\lambda_a\lambda_b},\quad a<b.
\end{equation}
Both verification scripts implement
Equation~\eqref{eq:curvature-diagnostic} directly.

\subsection{Empirical constant-curvature finding}\label{subsec:empirical-curv}

\begin{proposition}[Reproducible 4-bit curvature observation]\label{prop:empirical-curv}
For the PRESENT, GIFT, and SKINNY 4-bit S-boxes, the projected active-eigenspace
curvature diagnostic~\eqref{eq:curvature-diagnostic} gives
\begin{equation}\label{eq:empirical-minus-quarter}
  K_{ab}=-\frac14
\end{equation}
for all $\binom{15}{2}=105$ active coordinate planes, up to numerical precision
in the provided scripts.
\end{proposition}

\begin{proof}[Verification protocol]
The computation is reproducible from the repository scripts.  The scripts build
the WHT table, construct $\mathcal I^{\partial}$, diagonalize its active
15-dimensional eigenspace, project the third cumulant tensor, and evaluate
\eqref{eq:curvature-diagnostic}.  The reported value is stable across the three
tested S-boxes and an additional random 4-bit bijection in
\texttt{verify\_curvature.py}.

These reproducible computations suggest the following conjecture: for bijective
S-boxes, the active Fisher subspace associated with the Walsh boundary
covariance form has constant sectional curvature $K=-1/4$, under the curvature
convention and projection procedure used in this work.  A symbolic proof, or a
characterization of the precise hypotheses under which the conjecture holds, is
left as an open problem.
\end{proof}



\subsection{Relation to the belief simplex}\label{subsec:curvature-belief}

The belief simplex $Q=\Delta_{2^s-1}^{\circ}$ with Fisher--Rao metric is
isometric, via the square-root embedding, to a sphere of radius $2$ and hence
has constant sectional curvature $+1/4$.  The observed value $-1/4$ in the
active boundary covariance diagnostic suggests a sign-dual phenomenon between
the posterior belief geometry and the limiting keyed S-box covariance geometry,
which is recorded here as part of the open geometric structure of the model.
